% =====================================================================
\section{总结与思考}
% =====================================================================
\subsection{主要结论}
下面按五个研究问题逐一回答。所有结论都只针对本文的地图、发现模型和策略设置。

\textbf{RQ1（机制）。}在 A–G 楼这片校园里，躲确实比跑活得久：一轮 30 分钟里，就近躲进灌木或夹缝平均活 11.2 分钟，跑只活 5.4 分钟。反事实对照表明，这个优势主要来自“藏身处本身更难被看见”：只取消地形的隐蔽优势，躲的优势从 5.8 分钟降到 0.6 分钟；两种隐蔽优势都取消，就只剩 0.3 分钟，区间包含 0。只取消“站着不动更难被发现”这一项，躲仍多活 3.7 分钟。

\textbf{RQ2（最好的藏身点）。}就近选点时，夹缝最好（11.1 分钟，与“躲”相当），只躲灌木（6.3 分钟）反而不如原地不动（8.5 分钟）。但如果直接把老鼠分到藏身点，灌木和夹缝差不多一样安全（9.3 和 9.6 分钟，差别不显著）。所以好的藏身点要“藏得住、走得到”：夹缝紧贴每栋楼，走过去不用穿过大片空地。这与实战中鼠王躲在夹缝里、一直没被抓到的经验一致。

\textbf{RQ3（上下两轮之间的转化）。}不管第一轮有多少人躲，躲都比跑多活 4.5–5.5 分钟，是严格占优的策略。但如果每个人只看身边一个人的结果，两个方向的转化都会发生：第一轮四分之一的人躲时，15\% 的跑者改成躲，同时 34\% 的躲者改成跑。单局结果的运气成分很大，第二轮躲的比例只上升 0.01–0.02。

\textbf{RQ4（老鼠的学习）。}只奖励“活得久”的 AI 老鼠，部分学会了躲进夹缝、站着不动，但三个种子差别很大。它学到的是针对一起训练的猫的对策：换成脚本猫后，个人奖励组的 AI 老鼠只活 4.3 分钟，还不如脚本跑者。和历史版本的对手对局，比只和最新对手对局学得更通用（对上脚本猫 4.8 分钟对 2.3 分钟），但仍远不如脚本躲者。取消隐蔽优势后，AI 老鼠几乎不再进夹缝（18\% → 5\%），但冲刺没有增加，也没有变成“只会跑”。

\textbf{RQ5（猫的学习）。}团队平分奖励让猫明显变弱：脚本老鼠对上它能活 16.2 分钟，对上个人奖励的猫只活 8.6 分钟，三个种子无一例外。把功劳只分给抓捕时在场的猫，猫的强度就回到个人奖励的水平（9.3 分钟）；而把同样多的功劳随机分给几只猫，猫反而最弱（20.3 分钟）。所以起作用的不是奖励更集中，而是奖励落到了真正出力的猫身上，这支持“平分稀释了学习信号”的解释。不把围捕写进规则，AI 猫在抓捕时比脚本猫更聚集；团队奖励的猫最爱扎堆，抓得却最少。

\subsection{博弈论视角的思考}
这个游戏里有两层博弈。猫和老鼠之间，一开始是搜索博弈，一旦被发现就变成追逃博弈：“躲”把博弈尽量留在搜索阶段，“跑”则让自己更早进入追逃阶段。在体力有限、猫越来越多的条件下，追逃阶段对老鼠越来越不利，这就是“躲”占优的原因，而藏身处的隐蔽性决定了这个优势有多大。

老鼠和老鼠之间，被抓即转化在原理上带来外部性：每只老鼠被抓，其余老鼠就多一个敌人。但在本文的数值里，这种外部性很弱，每只老鼠的收益主要取决于自己选躲还是选跑，和别人怎么选关系不大。所以这里并不像公地悲剧那样“个人理性导致集体受损”：对个人最好的选择（躲）恰好也是对集体最好的选择。在本文的模仿模型里，真正阻碍“全员都躲”的是学习过程的局限：只看一局、只看一个人，运气会掩盖策略之间的真实差距。

猫的奖励设计可以类比机制设计中的激励结构：组织者通过奖励规则影响参与者的行为。不过，机制设计假设参与者是理性的，而本文的猫是靠强化学习去适应奖励的；团队奖励变弱的原因不是博弈论意义上的搭便车，而是学习信号被稀释。从博弈的角度看，团队奖励似乎更鼓励合作；但在学习意义上，奖励需要能追溯到个体的行为才有用。现实中“抓得最多的猫有奖”恰好给了每只猫清晰的激励。

\subsection{方法上的迭代与局限}
小场地阶段，AI 老鼠先后找到了两个环境漏洞：一是猫的巡逻点到不了墙角，老鼠就躲在角落里不动；二是猫只会直线追，老鼠绕着跑就能甩掉。给猫加上搜索记忆和预判拦截之后，结论才稳定下来。这说明对手太弱时，学到的“策略”可能只是漏洞。

本文还有几处局限。第一，发现概率的系数来自游戏体验，不是测量值；反事实对照说明了结论依赖哪一项，但具体的存活分钟数不能当真。第二，地图有一部分是合成的：OSM 是二维地图，没有空调外机的数据，空调外机是按固定随机种子沿外墙放置的；灌木丛来自 OSM 的林地和行道树，也不完全等于实地。第三，强化学习的规模有限：每组只有 3 个种子、300 轮（对手池两组 600 轮），同组种子之间差别很大，所以强化学习部分的比较都是探索性的。对手池组有两个模型训练中途中断过一次（一次是显卡驱动重置，一次是某一批只有一个有效样本、标准化时出现数值错误），都从最近的检查点“热重启”，环境和随机数状态没有恢复。第四，“按参与分配”只是局部共享奖励，不是精确衡量个体贡献的差分奖励；随机分配对照虽然支持信用稀释的解释，但每组只有 3 个种子，还不能算定论。第五，模拟本身还有两处近似：抓捕虽然在每个子步检查距离，但子步之间只按端点采样，斜着擦身而过的极近接触仍可能漏判；抓捕记给哪只猫、“按参与分配”的 15 m 范围和包围度，用的是这一步结束时的位置，而不是接触的那一刻，两者最多相差半秒。第六，两轮转化实验中的老鼠只有“躲”“跑”两种脚本策略，模仿规则也是假设的，真实的人还会听别人讲、会看整体结果。

下一步可以从三方面改进：用真正的差分奖励或值分解、更多的种子，确认团队奖励变弱的原因，并看能否在不削弱猫的情况下带来真正的围捕；扩大历史对手池、把脚本猫也放进对手池，看 AI 老鼠能不能学到对任何猫都管用的躲藏；回到实地标出真实的灌木与空调外机位置，再跑一次“哪里最好躲”。
